\chapter{线程池、任务切分与 NUMA}

单核内核写得再好，也只能使用一个核心。AI 推理要利用多核，必须把工作切成任务交给线程执行。但多线程不是免费加速。线程创建、任务分发、同步、缓存竞争、带宽饱和和 NUMA 跨节点访问都可能限制扩展。理解运行时，就是理解“哪些工作适合并行，怎样切，切到哪里，何时同步”。

\section{任务粒度}

任务太大，负载不均衡，一个线程可能长时间拖住整体；任务太小，调度和同步开销会压过计算。矩阵乘通常可以按 $M$ 或 $N$ 方向切块，卷积可以按 batch、输出通道或输出空间切块，Attention 可以按 head、batch 或 sequence block 切块。选择切分维度时，要看每个任务的工作量是否均衡，是否会重复读取大块数据，是否会让多个线程写同一 cache line。

线程池的基本思想是提前创建固定数量工作线程，避免每次算子调用都创建和销毁线程。任务队列保存待执行工作，工作线程取任务运行，主线程等待完成。高性能运行时还会做 work stealing、亲和性绑定、分层调度和 NUMA 感知分配。

\begin{lstlisting}[language=C++,caption={教学版并行 for 接口形状}]
#include <cstdint>
#include <functional>

class ParallelFor {
public:
    void run(std::int32_t begin,
             std::int32_t end,
             std::int32_t grain,
             const std::function<void(std::int32_t, std::int32_t)>& body);
};

void parallel_rows(ParallelFor& parallel_for,
                   std::int32_t rows,
                   const std::function<void(std::int32_t, std::int32_t)>& body) {
    constexpr std::int32_t ROW_GRAIN = 16;
    parallel_for.run(0, rows, ROW_GRAIN, body);
}
\end{lstlisting}

这只是接口形状，不是完整线程池实现。它强调一个重要点：并行库应该让调用者表达任务范围和粒度，而不是在算子内部随意创建线程。真实实现还需要异常传播、停止机制、队列、条件变量或无锁结构、线程生命周期和亲和性策略。

\section{带宽饱和}

有些内核单线程就接近内存带宽上限，多线程只会更快把带宽打满，超过某个线程数后不再加速。逐元素操作、拷贝、量化转换、某些 Softmax 和 KV Cache 扫描都可能出现这种情况。此时增加线程不但无益，还可能因为缓存争用和调度开销变慢。

判断带宽饱和要结合模型和测量。若每个元素只做很少计算，计算强度低，且多线程速度在某个线程数后平台化，就要怀疑带宽限制。可以用 STREAM 类 benchmark 估算机器内存带宽上限，再看内核实际带宽占比。若已经接近上限，下一步优化应减少字节数或融合算子，而不是继续增加线程。

\section{False Sharing 和写入冲突}

多个线程写不同变量，也可能互相影响。如果这些变量落在同一个 cache line 上，硬件缓存一致性协议会让 cache line 在核心之间来回转移，这就是 false sharing。并行写输出数组时，按连续大块切分通常比交错切分更安全。线程本地统计量也应避免紧挨在同一 cache line 上。

矩阵乘的并行切分还要避免多个线程同时更新同一块 $C$。若按 $K$ 维切分，每个线程计算部分和，最后需要归约，写冲突和额外内存都要处理；若按 $M$ 或 $N$ 块切分，每个线程负责独立输出块，通常更简单。不同形状下选择会变化，但原则是让写入所有权清晰。

\section{NUMA 的影响}

多路服务器或某些高端平台有 NUMA 结构：不同 CPU socket 或内存节点访问本地内存和远端内存的延迟、带宽不同。若线程在节点 0 上运行，却频繁访问节点 1 分配的内存，性能会下降。大模型权重和 KV Cache 占用巨大，NUMA 放置会影响端到端吞吐。

NUMA 优化包括首次触碰策略、线程绑定、按节点复制只读权重、按节点切分请求、避免跨节点频繁写共享数据等。是否值得做 NUMA 优化取决于部署环境。单手机或单 socket 机器上无需过早复杂化；多 socket 推理服务器上，NUMA 往往是必须解释的问题。

\begin{warningbox}
多线程扩展不是线性的。若一个内核从 1 到 4 线程加速明显，从 4 到 8 线程停滞，不要只怀疑线程池实现。先判断是否已经受内存带宽、LLC 容量、写冲突或 NUMA 限制。
\end{warningbox}

\begin{exercisebox}
对同一个向量加法、矩阵乘和 Softmax 分别测试 1、2、4、8 个线程。记录加速比，而不是只记录时间。解释哪个内核扩展最好，哪个最早平台化，并把原因写成计算强度和内存流量分析。
\end{exercisebox}
